\subsection{Which model is ``better''?}
On the decision metrics the ranking is unambiguous: Swin-T is the best single backbone
(\SwinAcc\% accuracy, \SwinBalAcc\% balanced accuracy, \SwinFone\% macro-F1), followed by
EfficientNet-B0 (\EffnetFone\%) and ResNet-50 (\ResnetFone\%). The gap between the transformer and
the best CNN is 6.3 macro-F1 points, which is larger than the 2.7 points separating the two CNNs,
so on this dataset the family label does matter---but not in the direction that is often assumed
for small medical datasets~\cite{matsoukas2021transformers}. Swin-T also converges faster
(\SwinFone\% macro-F1 is already reached at epoch 2 of 12 in validation, Fig.~\ref{fig:curves}),
which we attribute to its stronger ImageNet features and to the hierarchical windowed attention,
which lets it integrate colour and texture cues over the whole lesion rather than within a fixed
receptive field.

The ranking by macro AUC is the \emph{opposite} (ResNet-50 \ResnetAuc{} $>$ EfficientNet-B0
\EffnetAuc{} $>$ Swin-T \SwinAuc{}). AUC measures the quality of the probability ranking of each
class against the rest, whereas accuracy and F1 measure the hard arg-max decision. Swin-T fits the
training set much more tightly (its training loss ends at 1.00 versus 1.42 for ResNet-50 with the
same label smoothing), which produces sharper, more confident posteriors: better decisions on the
hard examples but slightly worse calibrated rankings on the easy ones. This is a reminder that a
single number can hide a trade-off, and that a clinical triage tool that thresholds probabilities
should be evaluated with AUC and calibration, not only with accuracy.

\subsection{Cost: FLOPs are not latency}
EfficientNet-B0 has 6$\times$ fewer parameters and 10$\times$ fewer FLOPs than ResNet-50, trains
in \EffnetTrain{} instead of \ResnetTrain{} minutes and uses the least VRAM
(\EffnetVram{}~GB). Yet on the GPU it is \emph{slower} at inference (\EffnetLat{}~ms versus
\ResnetLat{}~ms at batch~1; \EffnetThr{} versus \ResnetThr{} images/s at batch~32). Its
depthwise-separable convolutions are memory-bound and launch many small kernels that cannot
exploit tensor cores, so the theoretical savings do not materialise on this hardware; they would on
a CPU or mobile NPU, which is the setting EfficientNet was designed for~\cite{tan2019efficientnet}.
Swin-T has almost the same FLOPs as ResNet-50 (\SwinGflops{} versus \ResnetGflops{}~GFLOPs) but is
2.4$\times$ slower at batch~1, needs 1.5$\times$ more VRAM during training
(\SwinVram{} versus \ResnetVram{}~GB) and the longest wall-clock time (\SwinTrain{} min). The
accuracy gain of the transformer therefore has a concrete price: roughly 2.4$\times$ the latency and
1.5$\times$ the training memory of the residual CNN. In a batch setting (offline screening of a
clinic's daily images) this price is negligible; in an interactive or embedded setting it is the
deciding factor.

\subsection{Generalisation across classes}
All models are strong on the majority class NV (F1 0.85--0.88) and on the visually distinctive
VASC and BCC, and weak on the three rare or morphologically overlapping diagnoses AK, SCC and DF
(Table~\ref{tab:perclass}). The confusion matrices (Fig.~\ref{fig:cm}) show two persistent error
clusters: MEL$\leftrightarrow$NV (166 melanomas predicted as nevi and 159 nevi predicted as
melanoma by Swin-T) and the keratinocyte group AK$\leftrightarrow$BKL$\leftrightarrow$SCC, which is
also where dermatologists disagree most. Inverse-square-root class weighting raised the recall of
the rare classes (DF recall 0.60, VASC 0.81, SCC 0.46 for Swin-T) without collapsing the majority
class, but the balanced accuracy of 0.60--0.64 shows that the imbalance is only partly solved.
Swin-T is the most \emph{uniform} model: it has the best F1 in seven of the eight classes and is
the only backbone that lifts DF above 0.6, suggesting that its global receptive field helps on
classes defined by subtle, lesion-wide texture rather than by a local pattern.

\subsection{What the models look at}
The Grad-CAM maps (Fig.~\ref{fig:gradcam}) are consistent with that reading. The two CNNs produce
compact, centred blobs on the pigmented core of the lesion, which is adequate for NV, BCC, DF and
VASC but misses the periphery that distinguishes MEL from NV and AK from BKL. Swin-T produces
sparser, multi-focal maps (e.g. two separated hot spots on the melanoma, a ring on the vascular
lesion) that follow the irregular border and the asymmetric pigment distribution---precisely the
``ABCD'' cues used by clinicians. The failures are informative too: on the BKL example all three
models fire on the darkest region and predict melanoma, and on the SCC example EfficientNet-B0
spreads its evidence over the background, where hair and scale artefacts are present. Such
artefact-driven activations are a known risk on ISIC~\cite{cassidy2022isic} and argue for lesion
segmentation or artefact removal in a deployed system.

\subsection{Are the families complementary?}
The soft-voting ensemble improves on the best single model by 2.3 macro-F1 points
(\EnsFone\% versus \SwinFone\%), by 4.0 balanced-accuracy points and raises the macro AUC to
\EnsAuc{}, the best value in the study. Among the pairwise ensembles, the two CNNs together reach
only 61.5\% macro-F1, whereas EfficientNet-B0+Swin-T reaches 66.4\%, almost the value of the full
ensemble; i.e. nearly all the gain comes from combining a CNN with the transformer. The two
families make different mistakes, which is the strongest practical argument for hybrid systems,
at the cost of \EnsLat{}~ms per image.

\subsection{Comparison with the literature and limitations}
Our balanced accuracies (0.60--0.68) are in the range reported for single ImageNet-pretrained
models on ISIC~2019 at 224~px, and below the 0.63--0.74 obtained by the challenge winner with
multi-resolution EfficientNet ensembles trained on far more compute and with
metadata~\cite{gessert2020isic2019}; note that our numbers are measured on an internal,
in-distribution split and are therefore not directly comparable to the official out-of-distribution
test set. Limitations: a single seed per model (differences below $\approx$1 point should not be
over-interpreted), a fixed 224~px resolution that discards fine dermoscopic structure, no use of
age/sex/site metadata, no test-time augmentation, and a learning rate that was not tuned per
architecture (a per-model search would most likely narrow, not widen, the gaps).
